\subsection{Integration with respect to an induced measure}

Let X be a locally compact subspace of T, \(\mu\) a positive measure on T, and \(\mu_{X}\) the measure induced on X by \(\mu\) (Ch. IV, §5, No.~7). For every \(t \in T\) , let us define a measure \(\lambda_{t}\) on X in the following way: \(\lambda_{t} = \varepsilon_{t}\) if \(t \in X\) , \(\lambda_{t} = 0\) if \(t \in C X\) . For every finite numerical function g defined on X, \(\int g(x) d\lambda_{t}(x) = g(t)\) if \(t \in X\) and \(\int g(x) d\lambda_{t}(x) = 0\) if \(t \in C X\) . If g is a function in \(\mathcal{K}(X)\) we therefore have, by the definition of \(\mu_{X}\) ,

\[
\mu_ {\mathrm{X}} (g) = \int \langle g, \lambda_ {t} \rangle d \mu (t).\tag{1}
\]

This means that one can write

\[
\mu_ {\mathrm{X}} = \int \lambda_ {t} d \mu (t)\tag{2}
\]

(§3, No.~1).

Let us now define a mapping \(\pi\) of T into X by setting \(\pi(t) = t\) for \(t \in \mathbf{X}\) , and \(\pi(t) = t_0\) for \(t \in \mathbb{C} \mathbf{X}\) , \(t_0\) being an arbitrary point of \(\mathbf{X}\) ; one can write \(\lambda_t = \varphi_{\mathbf{X}}(t) \varepsilon_{\pi(t)}\) for every \(t \in \mathbf{T}\) . The mapping \(\pi\) is \(\mu\) -measurable, because its restrictions to \(\mathbf{X}\) and \(\mathbb{C} \mathbf{X}\) are (Ch. IV, §5, No.~10, Prop. 16); it follows at once that the pair \((\pi, \varphi_{\mathbf{X}})\) is \(\mu\) -adapted (§4, No.~1). We therefore have the following results:

PROPOSITION 1. --- For every numerical function \(g \geqslant 0\) defined on X,

\[
\int^ {\bullet} g d \mu_ {\mathrm{X}} = \int_ {\mathrm{X}} ^ {\bullet} g d \mu\tag{3}
\]

(cf.~§5, No.~3, Example, for the notation \(\int_{\mathbf{X}}^{\bullet}\) ).

Taking into account the preceding remarks and (2), the relation (3) follows from Th. 1 of §4.

COROLLARY 1. --- For every subset B of X, \(\mu_{\mathrm{X}}^{\bullet}(\mathrm{B}) = \mu^{\bullet}(\mathrm{B})\) ; for B to be locally \(\mu_{\mathrm{X}}\) -negligible, it is necessary and sufficient that B be locally \(\mu\) -negligible.

COROLLARY 2. --- Let M be a subset of T. If \(\mu\) is concentrated on M, then \(\mu_{X}\) is concentrated on \(M \cap X\) .

COROLLARY 3. --- For the measure \(\mu_{X}\) to be zero, it is necessary and sufficient that \(X\) be locally \(\mu\) -negligible.

Remark. --- If S is the support of \(\mu\) , then \(S \cap X\) (which is closed in X) contains the support of \(\mu_{X}\) by Cor. 2, but may be distinct from it. For example, if \(\mu\) is a diffuse measure and X is a subspace reduced to a point, then the induced measure \(\mu_{X}\) is zero, hence its support is empty. Note, however, that the support of \(\mu_{X}\) is equal to \(S \cap X\) if X is open in T.

PROPOSITION 2. --- For a mapping g of X into a topological space to be \(\mu_{X}\) -measurable, it is necessary and sufficient that g be \(\mu\) -measurable in X ( \(\S5\) , No.~3, Example).

This follows from Prop. 3 of §4.

COROLLARY. --- For a subset B of X to be \(\mu_{X}\) -measurable, it is necessary and sufficient that B be \(\mu\) -measurable.

THEOREM 1. --- Let g be a function defined on X, with values in \(\overline{R}\) or in a Banach space. For g to be essentially \(\mu_{X}\) -integrable, it is necessary and sufficient that \(\mathbf{g}\) be essentially \(\mu\) -integrable in X (§5, No.~3, Example), in which case

\[
\int \mathbf {g} d \mu_ {\mathrm{X}} = \int_ {\mathrm{X}} \mathbf {g} d \mu .\tag{4}
\]

This follows from Th. 2 of §4.

COROLLARY 1. --- For a subset B of X to be essentially \(\mu_{X}\) -integrable, it is necessary and sufficient that it be essentially \(\mu\) -integrable, in which case \(\mu_{\mathrm{X}}(\mathrm{B}) = \mu(\mathrm{B})\) .

COROLLARY 2. --- Let g be a complex function defined on T and locally \(\mu\) -integrable; the restriction \(g_{X}\) of g to X is then locally \(\mu_{X}\) -integrable, and

\[
(g \cdot \mu) _ {\mathrm{X}} = g _ {\mathrm{X}} \cdot \mu_ {\mathrm{X}}.\tag{5}
\]

This follows at once from Th. 1, applied to the functions \(fg_{\mathbf{X}}\) ( \(f \in \mathcal{K}(\mathbf{X}; \mathbf{C})\) ), and the definition of the measure induced on \(\mathbf{X}\) by a complex measure (Ch. IV, §5, No.~7).

COROLLARY 3. --- Let \(\theta\) be a complex measure on T; then

\[
| \theta | _ {\mathrm{X}} = | \theta_ {\mathrm{X}} |.\tag{6}
\]

Set \(|\theta| = \mu\) and apply Cor. 2 on taking \(g\) to be a complex function of absolute value 1 such that \(\theta = g \cdot \mu\) (§5, No.~5, Cor. 3 of Th. 2); then \(\theta_{\mathrm{X}} = g_{\mathrm{X}} \cdot \mu_{\mathrm{X}}\) ; but \(g_{\mathrm{X}}\) is a function of absolute value 1, and the formula (6) follows from Prop. 2 of §5, No.~2.

Remarks. --- a) Cor. 3 has already been proved by another method (Ch. IV, §5, No.~7, Lemma 3).

\begin{enumerate}
\def\labelenumi{\alph{enumi})}
\setcounter{enumi}{1}
\tightlist
\item
  By virtue of Cor. 3, the corollaries 1,2,3 of Prop. 1, Prop. 2, Theorem 1 and its corollaries 1,2 extend at once to a complex measure.
\end{enumerate}

Scholium. --- For every function \(\mathbf{f}\) (resp. \(\mathbf{g}\) ) defined on X (resp. T) with values in the Banach space F or in \(\overline{\mathbf{R}}\) , let us denote by \(\zeta(\mathbf{f})\) (resp. \(\rho(\mathbf{g})\) ) the extension by 0 of \(\mathbf{f}\) to T (resp. the restriction of \(\mathbf{g}\) to X). Then \(\zeta(\rho(\mathbf{g})) = \varphi_{\mathrm{X}} \cdot \mathbf{g}\) , \(\rho(\zeta(\mathbf{f})) = \mathbf{f}\) . We denote by \(\mu'\) the measure \(\varphi_{\mathrm{X}} \cdot \mu\) on T. For every \(p \in [1, +\infty]\) , Props. 1 and 2 imply that \(\zeta\) maps \(\overline{\mathcal{L}}_{\mathrm{F}}^{p}(\mathrm{X}, \mu_{\mathrm{X}})\) into \(\overline{\mathcal{L}}_{\mathrm{F}}^{p}(\mathrm{T}, \mu')\) , and that \(\rho\) maps \(\overline{\mathcal{L}}_{\mathrm{F}}^{p}(\mathrm{T}, \mu')\) onto \(\overline{\mathcal{L}}_{\mathrm{F}}^{p}(\mathrm{X}, \mu_{\mathrm{X}})\) , with preservation of norm in both cases, as well as of the integral when \(p = 1\) (Th. 1); passing to the associated Hausdorff spaces, we obtain two isomorphisms inverse to each other. Similarly, if \(\zeta\) and \(\rho\) are applied to positive numerical functions, the essential upper integral is preserved (Prop. 1). Thus, if we agree to identify a function on X with a function on T that is zero on X-T, and the measure \(\mu_{X}\) with the measure \(\mu'\) , problems concerning induced measures are reduced to problems concerning measures defined by densities, treated in §5. This sort of reasoning is applicable to complex measures as well, by Cor. 3 of Th. 1.
% semantic-fragments: legacy-input-seam
\relax{}
